\documentclass[11pt]{article}
\usepackage[a4paper,margin=25mm]{geometry}
\usepackage{amsmath,amssymb,amsthm}
\usepackage[hidelinks]{hyperref}
\emergencystretch=2em
\newtheorem{theorem}{Theorem}
\title{Random weighted cycle tree polynomials\newline
\large Disproof of Conjecture 00000000470}
\author{gaochengzhecpu}\date{}
\begin{document}\maketitle
\begin{abstract}
Give every edge of the cycle $C_m$ an independent fair weight
in $\{1,2\}$. With probability $1-2^{1-m}$ its weighted
spanning-tree polynomial has exactly one nonzero root, a
strictly negative real number. The nonzero-root angular
empirical measure is therefore a point mass at $\pi$, at
Kolmogorov distance at least $1/2$ from uniform measure.
The same fixed gap holds if zero roots are included with
argument zero. Both genuine nonempty empirical measures
fail with probability tending to one. Thus even convergence
in probability to uniformity fails for this vertex-transitive
family, excluding the proposed vanishing rate.
\end{abstract}

\section*{The graph family and the independent weights}
For every integer $m\geq3$, take the simple cycle with vertex
set $\mathbb Z/m\mathbb Z$ and undirected edges
$e_i=\{i,i+1\}$. The rotations $i\mapsto i+d$ are graph
automorphisms and carry any specified vertex to any other.
These are genuine vertex-transitive graphs, with $m$ vertices
and $m$ distinct edges, including when $m=3$.

Assign weight $w_i\in\{1,2\}$ independently and uniformly
to each original edge. More explicitly, the sample space is
$\{1,2\}^{m}$, every complete assignment having probability
$2^{-m}=\prod_{i=0}^{m-1}(1/2)$. This is a nondegenerate
identically distributed independent law; no conditioning is
part of the model. Write
\[
 a=\#\{i:w_i=1\},\qquad b=\#\{i:w_i=2\},\qquad
 W=\sum_i w_i=a+2b.
\]
Then $a+b=m$ and $W\geq m$. The mixed event
$M_m=\{a>0,\ b>0\}$ excludes exactly the two constant
assignments, so
\[
                 \Pr(M_m)=1-\frac{2}{2^m}\longrightarrow1.
\]
The source places no additional restriction on the edge-weight
law that would exclude these positive integer weights.

\newpage
\section*{The actual spanning-tree polynomial and its roots}
Every spanning tree of a cycle is obtained by deleting
exactly one edge. Indeed, deleting any edge leaves a connected
path with $m-1$ edges, hence a tree. Conversely a spanning tree
has $m-1$ edges, and therefore omits exactly one of the
cycle's $m$ edges. Distinct omitted edges give distinct trees.
Consequently the polynomial defined in the source is exactly
\begin{align*}
 \tau_p(C_m)&=\sum_{T\text{ spanning tree}}p^{\sum_{e\in T}w_e}
          =\sum_{i=0}^{m-1}p^{W-w_i}\\
          &=a p^{W-1}+b p^{W-2}
           =p^{W-2}(ap+b).
\end{align*}
This calculation sums all actual spanning trees, so coincident
exponents accumulate with their full multiplicities.

On $M_m$, both $a$ and $b$ are positive. The polynomial is
nonzero, of degree $W-1$, and its complete root multiset is
\[
        \underbrace{0,\ldots,0}_{W-2\text{ copies}},
                  \quad -\frac ba.
\]
The last root is simple, nonzero and strictly negative real.
Its argument is $\pi$ in both the usual $[0,2\pi)$ convention
and the principal argument convention. Thus no choice between
these usual branches affects the example.

\section*{A fixed Kolmogorov gap under both zero conventions}
For a nonempty multiset $Z$ of roots with chosen arguments,
counted with multiplicities, write
\[
 F_Z(x)=\frac{\#\{z\in Z:\arg z\leq x\}}{|Z|},\qquad
 U(x)=\begin{cases}0&x<0,\\x/(2\pi)&0\leq x\leq2\pi,\\1&x>2\pi.
 \end{cases}
\]
The Kolmogorov distance is the full supremum
$D(Z)=\sup_{x\in\mathbb R}|F_Z(x)-U(x)|$.
It always dominates the error at any individual threshold.

If zero roots are discarded, $Z=\{-b/a\}$ on $M_m$.
It is nonempty, and $F_Z(\pi)=1$ whereas $U(\pi)=1/2$.
Thus $D(Z)\geq1/2$.

Alternatively, include every root and explicitly set
$\arg(0)=0$. The complete multiset is nonempty, and
\[
 F_Z(0)=\frac{W-2}{W-1}\geq\frac12,\qquad U(0)=0,
\]
because $W\geq m\geq3$. Again $D(Z)\geq1/2$.
The first argument already gives a disproof without assigning
any argument to zero.

\newpage
\section*{The probabilistic obstruction and the claimed rate}
For either convention, let $E_m$ be the event that its actual
root multiset is nonempty and its Kolmogorov distance is at
least $1/2$. The preceding computation proves $M_m\subseteq E_m$,
and hence
\[
              1-2^{1-m}\leq\Pr(E_m)\leq1,
                  \qquad \Pr(E_m)\longrightarrow1.
\]
In particular these probabilities do not tend to zero.
The positive-cardinality condition ensures that every
empirical measure used in this conclusion is defined in
the usual sense. No empty multiset or totalized division
can supply the obstruction.

\begin{theorem}
For the vertex-transitive family $C_m$, $m\geq3$, with
independent fair weights in $\{1,2\}$, the empirical root
arguments fail to approach the uniform distribution even
in probability. This holds after discarding zero roots,
and also when retaining them with argument zero.
\end{theorem}
\begin{proof}
For either convention a fixed Kolmogorov gap of at least
$1/2$ occurs on the mixed event, whose probability tends
to one. This contradicts convergence in probability to
the uniform distribution.
\end{proof}
For any fixed finite constant $C\geq0$, eventually
$C m^{-1/2}<1/2$. The same mixed events then violate a
bound $D\leq C m^{-1/2}$. Thus the asserted vanishing
$O(m^{-1/2})$ rate cannot repair the failure. The formal
proof establishes the stronger fixed-gap probability
limit; this elementary comparison with $m^{-1/2}$ is a
consequence stated here, not a separate Lean theorem.

The construction refutes the conjecture with its supplied
definition $\tau_p(G)=\sum_T p^{w(T)}$. It makes no claim
about a different polynomial or a separately specified
weight law. Real coefficients make roots conjugation
invariant; here all roots in the decisive event are real.

\section*{Formalization and checks}
The self-contained Lean project uses Mathlib's actual simple
cycle graph on $\mathbb Z/(n+3)\mathbb Z$. It constructs
rotation graph isomorphisms and proves that all genuine
spanning-tree subgraphs are bijective with deleted original
edges. Tree weights sum over the actual undirected edge
sets. The polynomial is the sum over all those actual trees;
its factorization and complete complex root multiset are
derived, rather than assumed.

\newpage
\section*{Formalization and reproducibility, continued}
The formal empirical CDF uses the actual root multiset,
its filtered cardinality at each threshold, and its full
cardinality. The Kolmogorov distance is a supremum over
all real thresholds; the range is proved bounded above
before using the supremum lower bound. The source's
$[0,2\pi)$ convention and Mathlib's principal convention
agree on every root at issue.

Normalized counting on all Boolean edge assignments is
the finite uniform product law. Lean proves total mass one,
probability bounds, monotonicity, and that the probability
of every complete assignment is the product of its fair
coordinate masses. It then proves the exact mixed-event
count, its probability limit, and for each zero convention
the probability limit one for the actual nonempty
fixed-gap event. Uniqueness of limits proves that these
probabilities cannot tend to zero.

The accompanying \texttt{verify.py} independently enumerates
actual connected $(m-1)$-edge subgraphs and every weight
assignment for $3\leq m\leq10$. Exact integer and rational
checks confirm the spanning-tree enumeration, polynomial
coefficients, root description, CDF lower bounds and mixed
event probability. These finite checks are supplementary;
Lean proves the all-size result.

Lean 4.19.0 and Mathlib revision\\
\texttt{c44e0c8ee63ca166450922a373c7409c5d26b00b}
are pinned. From \texttt{lean/}, run
\begin{verbatim}
lake exe cache get
lake build
lake env lean -DwarningAsError=true Main.lean
\end{verbatim}
The project is rebuilt afresh, and the printed theorem
dependencies contain only standard foundational axioms.
There are no proof placeholders, custom axioms, external
computation oracles or imports from another submission.
Compile this paper with \texttt{tectonic main.tex}.

Exact bilingual source, provenance, real verification logs
and artifact hashes are included. The author agent and
parent performed separate internal reviews; no external
independent review is claimed. The elementary cycle-cut
lemmas were copied and adapted from our earlier 7841
submission, with their complete proofs included here.

\begin{thebibliography}{9}
\bibitem{source} The Last Math Competition,
Conjecture 00000000470. Exact bilingual source included
as \texttt{SOURCE.md}. All facts needed for this example
are proved in this paper and the accompanying Lean project.
\end{thebibliography}
\end{document}
